\section{Subject-Specific Equivalent Dynamics Identification}
\label{sec:identification}

\subsection{Identification Data and Excitation Design}
\label{subsec:identification_data}

% Trace: IV-A-P1
Each identification sample contains the joint state and the sagittal-plane
interaction force,
\begin{equation}
    \vect{y}_i=
    \left\{\vect{q}_i,\dot{\vect{q}}_i,\ddot{\vect{q}}_i,
    \vect{F}_{R\rightarrow L,i},t_i,v_i\right\},
    \label{eq:identification_signals}
\end{equation}
where $v_i$ is an explicit validity flag.  Invalid, non-finite,
near-singular, or stale observations are excluded rather than replaced by
zeros.  The measured generalized torque used for fitting is always rebuilt as
$\vect{\tau}_{\mathrm{meas},i}=\vect{J}(\vect{q}_i)^{\mathsf T}
\vect{F}_{R\rightarrow L,i}$; stored generator torques and virtual-subject
parameters are not estimator inputs.

% Trace: IV-A-P2
The repository contains two complementary excitation designs.  The general
software identification suite combines three joint-motion families
(coupled, hip-dominant, and knee-dominant) with slow, nominal, and fast timing.
The nine complete trajectories are assigned before fitting: four trajectories
to training, two to validation, and three to testing.  This variation in joint
amplitude, relative hip--knee motion, velocity, and acceleration is intended to
separate inertial, stiffness, and damping effects while remaining within the
implemented workspace and ROM.

% Trace: IV-A-P3
The current task-local suite is centered on the frozen active asymmetric
reference and contains 12 prespecified 401-sample trajectories.  Six training
trajectories apply endpoint-zero $3^{\circ}$ hip- or knee-amplitude reductions
or a branch-local monotone $\pm3\%$ knee-phase warp at slow or nominal timing.
Two intermediate trajectories, with a $2^{\circ}$ hip reduction or $+2\%$
knee-phase warp, form the validation split.  Testing uses an unseen $-2\%$
knee-phase warp, the exact active slow reference, its exact relevant nominal
profile, and a nominal profile retimed $10\%$ faster.  Neither exact reference
profile enters parameter fitting.  The state-domain gate is an axis-aligned
box in $[q_h,q_k,\dot q_h,\dot q_k,\ddot q_h,\ddot q_k]$ fitted from training
states only.

% Trace: IV-A-P4
These two suites answer different questions.  The general suite supports
identifiability and robustness stress tests, whereas the current local suite
tests train-only transfer near the prescribed active task.  All local
trajectories retain the active-reference identifier and file hash, and the
legacy symmetric local dataset is excluded from fitting.  Its stored results
are used only for a retrospective comparison under common metric definitions;
they are not evidence that asymmetric trajectory geometry is intrinsically
superior.

\subsection{Equivalent Parameter Estimation}
\label{subsec:parameter_estimation}

% Trace: IV-B-P1
The estimated parameter vector is
\begin{equation}
    \vect{\vartheta}=
    \begin{bmatrix}
    \beta_m & k_h & k_k & b_h & b_k
    \end{bmatrix}^{\mathsf T},
    \label{eq:parameter_vector}
\end{equation}
where $\beta_m$ jointly scales both segment masses and both planar inertias of
a fixed baseline template.  Center-of-mass distances, neutral angles, and
gravity remain fixed.  Given an observed state, the predicted torque is
\begin{equation}
    \hat{\vect{\tau}}_i(\vect{\vartheta})
    =\vect{f}_{\mathrm{dyn}}
    (\vect{q}_i,\dot{\vect{q}}_i,\ddot{\vect{q}}_i;
     \vect{\vartheta}),
    \label{eq:torque_prediction}
\end{equation}
where $\vect{f}_{\mathrm{dyn}}$ is exactly the inverse-dynamics model in
Section~\ref{subsec:dynamics}.

% Trace: IV-B-P2
Only training observations are passed to the optimizer.  Let
$\sigma_h=\max(\operatorname{std}(\tau_{h,\mathrm{meas}}),
\SI{1}{N.m})$ and define $\sigma_k$ analogously.  The implemented estimate is
\begin{equation}
    \hat{\vect{\vartheta}}
    =\arg\min_{\vect{\vartheta}\in\Theta}
    \sum_{i\in\mathcal{D}_{\mathrm{train}}}
    \left[
    \rho_{\mathrm{SL1}}\!\left(
      \frac{\tau_{h,i}-\hat\tau_{h,i}}{\sigma_h}
    \right)
    +
    \rho_{\mathrm{SL1}}\!\left(
      \frac{\tau_{k,i}-\hat\tau_{k,i}}{\sigma_k}
    \right)
    \right],
    \label{eq:estimation_objective}
\end{equation}
where $\rho_{\mathrm{SL1}}$ denotes the implemented soft-$\ell_1$ robust loss.
The joint-specific scaling prevents one torque channel from dominating solely
because of its numerical range.

% Trace: IV-B-P3
The parameter set is the box
\begin{equation}
    \Theta=[0.6,1.6]
    \times[0,60]^2\ \si{N.m.rad^{-1}}
    \times[0,10]^2\ \si{N.m.s.rad^{-1}}.
    \label{eq:parameter_bounds}
\end{equation}
The initial vector is
$[1,10,10,1,1]^{\mathsf T}$ in the corresponding units, and the numerical
solver uses parameter scaling $[1,20,20,3,3]^{\mathsf T}$ and at most 500
function evaluations.  Validation and test samples are used only after the
train-only fit to compute RMSE, MAE, $R^2$, and variance-accounted-for metrics.

% Trace: IV-B-P4
This estimator is a bounded nonlinear least-squares gray-box method.  The
current implementation does not use a physics-informed neural network,
Bayesian estimator, comfort predictor, or model-predictive controller.  The
reported parameters must therefore be interpreted as task-local equivalent
coefficients, particularly when the data-generating mechanics violate the
five-parameter structure.

\subsection{Parameter Identifiability}
\label{subsec:identifiability}

% Trace: IV-C-P1
Local identifiability is evaluated from a numerical sensitivity matrix.  For
the normalized residual vector $\vect{r}$ and parameter scaling matrix
$\vect{S}_{\vartheta}$, the implemented analysis approximates
\begin{equation}
    \vect{H}=
    \frac{\partial\vect{r}}{\partial\vect{\vartheta}}
    \vect{S}_{\vartheta}
    \label{eq:sensitivity_matrix}
\end{equation}
by central finite differences when both perturbation directions satisfy the
bounds, and by a one-sided difference otherwise.  The same fixed hip and knee
torque scales are used when comparing excitation subsets.

% Trace: IV-C-P2
The singular values of $\vect{H}$, its numerical rank, and
$\kappa(\vect{H})$ quantify local parameter-direction sensitivity.  The
information-shape matrix and its pseudoinverse are
\begin{equation}
    \vect{I}=\vect{H}^{\mathsf T}\vect{H},
    \qquad
    \vect{C}_{\vartheta}\propto\vect{I}^{\dagger},
    \label{eq:information_matrix}
\end{equation}
from which normalized parameter correlations are computed.  The software
compares a single coupled nominal trajectory, the coupled family at all three
speeds, all nine trajectories, and the latter set after removing the most
extreme $5\%$ of Jacobian condition numbers.  It also reports sensitivity by
observed-force quartile without deleting high-force observations.

% Trace: IV-C-P3
A low training or test prediction error is not sufficient evidence of unique
physical parameter recovery.  Correlated columns of $\vect{H}$ can produce
similar torque predictions for different parameter combinations, and the
local covariance approximation is conditional on the selected model and
excitation.  Rank, conditioning, correlation, held-out prediction, and
parameter-bound activity must therefore be reported together.  None of these
software diagnostics turns an equivalent coefficient into physiological
ground truth.

\subsection{Timing Alignment and Delay Robustness}
\label{subsec:timing}

% Trace: IV-D-P1
State and wrench streams are treated as asynchronous.  For the fixed-delay
study, positive delay is defined by
\begin{equation}
    \vect{F}_{\mathrm{obs}}(t)
    =\vect{F}_{\mathrm{clean}}(t-\delta).
    \label{eq:fixed_delay_definition}
\end{equation}
Using the wrong $\delta$ assigns a force generated by one state to another
state, creating a phase-dependent torque residual.  The repository searches a
$1$-ms grid from $-50$ to $50$ ms.  For each candidate, the five parameters are
fit using training data, the candidate is selected by validation torque RMSE,
and the test split remains unavailable until selection is complete.

% Trace: IV-D-P2
All fixed-delay candidates are compared on common valid support.  A
\SI{50}{ms} boundary margin is used, common-support coverage must be at least
$80\%$, and offline interpolation does not bridge a gap longer than
\SI{20}{ms}.  A selected grid endpoint is flagged as a possible out-of-range
delay rather than accepted as an exact estimate.  The offline interpolation
mode may use a future-arriving sample and is therefore restricted to offline
identification; a separate causal-history mode uses only already available
samples.

% Trace: IV-D-P3
The variable-delay study distinguishes state source time, wrench sample time,
and wrench arrival time.  When a reliable wrench sample timestamp is present,
the target state time is that sample time; otherwise it is estimated as
\begin{equation}
    t_{q}^{\star}=t_{F,\mathrm{arr}}-\hat\delta(t).
    \label{eq:variable_delay_target}
\end{equation}
A \SI{2}{s} state-history buffer permits nearest-neighbor or bracketed linear
matching but rejects future queries, expired history, interpolation gaps over
\SI{20}{ms}, and state-match errors over \SI{5}{ms}.  Dropout, stale freeze,
excess wrench age, duplicate sample timestamps, and non-causal timestamps have
explicit invalid reasons.

% Trace: IV-D-P4
The windowed delay tracker uses a \SI{2}{s} window, updates every
\SI{0.5}{s}, searches $[-50,80]$ ms at 1-ms resolution, applies exponential
smoothing with coefficient $0.5$, and limits an accepted update to
\SI{8}{ms}.  It requires at least 25 effective samples and sufficient velocity
and acceleration excitation; otherwise the previous estimate is held with an
invalid-update reason.  These fixed- and variable-delay experiments serve only
to characterize the sensitivity of equivalent identification to temporal
misalignment.  They are not the headline contribution and do not establish
real-time robot performance.

\subsection{Model-Mismatch and Generalization Analysis}
\label{subsec:model_mismatch}

% Trace: IV-E-P1
To test whether a compact equivalent model remains predictive when its
assumptions are violated, the virtual data generator augments the linear model
with selected terms:
\begin{equation}
    \vect{\tau}_{\mathrm{gen}}
    =\vect{\tau}_{5\mathrm{p}}
    +\vect{\tau}_{K_3}
    +\vect{\tau}_{K_c}
    +\vect{\tau}_{B_2}
    +\vect{\tau}_{r}.
    \label{eq:mismatch_generator}
\end{equation}
Here, $\vect{\tau}_{K_3}$ contains independent cubic stiffness terms,
$\vect{\tau}_{K_c}$ is derived from a hip--knee coupling potential,
$\vect{\tau}_{B_2}$ contains odd-symmetric quadratic damping, and
$\vect{\tau}_{r}$ is a smooth deterministic structured residual.  The
estimator remains the five-parameter model in~\eqref{eq:parameter_vector}; none
of the additional generator terms, scenario labels, or generator parameters is
available to it.

% Trace: IV-E-P2
The implemented suite contains a matched-linear case; mild and strong
nonlinear stiffness; mild and strong hip--knee coupling; mild nonlinear
damping; a structured residual; and mild and strong combined mismatch.  Only
the prescribed training split is fitted.  Evaluation includes validation plus
unseen interpolation paths with phase, amplitude, speed, and asymmetric
flexion--extension changes, a near-boundary path, and a deliberately labeled
outside-domain path.

% Trace: IV-E-P3
The identified equivalent model and a fixed generic baseline are evaluated on
the same valid samples using torque RMSE, endpoint-force RMSE, peak error, and
normalized RMSE.  The outside-domain set is retained to expose extrapolation
risk and is not used to assert task-local validity.  When residuals correlate
with angle, speed, or the opposite-joint state, the implementation records a
possible structural-mismatch diagnostic rather than inferring a physiological
mechanism.

% Trace: IV-E-P4
Under mismatch, $\hat{\vect{\vartheta}}$ is the best equivalent predictive
model over the training region, not a recovery of the generator's nonlinear
coefficients or true human physiology.  Useful held-out prediction can coexist
with biased individual coefficients, whereas strong structured mismatch can
make the identified model worse than the generic baseline on some trajectories.
The generalization analysis must therefore report failure cases and domain
membership together with average improvements.
